\documentclass[11pt]{article}

\usepackage[margin=1in]{geometry}
\usepackage{graphicx}
\usepackage{booktabs}
\usepackage{amsmath}
\usepackage{hyperref}
\usepackage{siunitx}
\usepackage{float}
\usepackage{xurl}

\hypersetup{colorlinks=true, linkcolor=blue, citecolor=blue, urlcolor=blue}

\title{Physically Motivated Time-Series Photometry of AM CVn}
\author{TODO: Team Members}
\date{TODO: Submission Date}

\begin{document}
\maketitle

\begin{abstract}
We observed AM Canum Venaticorum (AM CVn) in the V band on 2026 May 20 to test
whether short-baseline CCD photometry can recover and refine its known
ultracompact variability.  We reduced 140 images with bias, dark, and
star-masked flat-field calibration, then performed AAVSO-sequence differential
photometry.  After rejecting two artifact frames, 138 exposures spanning
111.20 min remained.  A generic single-sine search detects a 525.06 s signal,
close to the expected first harmonic of AM CVn's positive superhump.  A
physically motivated double-wave model gives the best independent period,
\(1055.21^{+13.10}_{-9.63}\) s from fixed-time Monte Carlo perturbations after
empirical error scaling.  After penalizing
free/scanned periods with an effective BIC, a fixed 1051.2 s
superhump-plus-harmonic template is the best conditional model, while the free
double wave remains the best independent period measurement.  The data
therefore support detection of AM CVn's known superhump/harmonic family, but
they do not tighten the literature period constraints.
\end{abstract}

\section{Introduction}

AM CVn is the prototype of a class of hydrogen-deficient ultracompact
interacting binaries.  These systems contain an accreting white dwarf and a
helium-rich donor, and their periods are short enough to connect accretion-disc
physics, compact-binary evolution, and low-frequency gravitational-wave
astrophysics \cite{solheim2010,roelofs2006}.

The period structure of AM CVn is subtle.  Early photometry showed a double
wave near 17.5 min and a strong 8.76 min harmonic \cite{patterson1979}.  The
modern period structure was laid out by Skillman et al., who reported the
positive superhump at 1051.2 s, its strong 525.6 s harmonic, and the precise
1028.7322 s clock \cite{skillman1999}.  Spectroscopy then established the
1028.73 s signal as the true orbital period, while the dominant 1051.2 s
photometric signal is interpreted as a positive superhump
\cite{nelemans2001,roelofs2006}.  Recent work uses this period structure to
revisit AM CVn system parameters and period change \cite{smak2023}.

Our goal is not only to detect variability, but to ask a sharper question: can
our local 111 min light curve constrain the literature periods more tightly?
To avoid overclaiming, we separate independent constraints from conditional
constraints.  Independent constraints use periods fitted freely from our data.
Conditional constraints assume a literature-motivated physical identification,
such as the known superhump or orbital signal.

\section{Observations and Reduction}

The observations were taken on 2026 May 20 in the V band.  The AM CVn sequence
contains 140 frames with exposure times of 35, 40, and 45 s.  The reduction used
late-night bias frames, exposure-matched darks, and a star-masked
\texttt{flat\_4s\_masked\_clipmean} flat.  The final two frames showed severe
artifacts and low counts, so they were excluded.

The science frames did not have a reliable WCS solution.  We therefore used the
AAVSO comparison sequence X42421BZ to map the target and comparison stars onto
a good reference frame.  Aperture photometry used a 6 px aperture and a
12--20 px sky annulus.  Fluxes were divided by exposure time before comparison,
so frames with different exposure times could be combined.  Each comparison
star was normalized by its own median flux before constructing the ensemble,
preventing the brightest comparison stars from dominating the differential
light curve.

The final light curve contains 138 good points.  It spans 111.20 min, or 6.35
cycles of a 1051.2 s superhump.  The robust target scatter is 1.51 percent,
while the median formal photometric error is 0.90 percent.  This excess scatter
motivates the empirical error scaling used below.

As a control on the differential photometry, we also treated each comparison
star as a pseudo-target and divided it by the remaining comparison ensemble.
Those pseudo-light curves are consistent with constant-flux controls at the
level needed for this project: the six comparison stars have robust scatters of
0.70--1.22 percent and empirical constant-model reduced chi-squared values near
unity.

\begin{figure}[H]
\centering
\includegraphics[width=\textwidth]{analysis/am_cvn/plots/AM_CVn_comparison_star_controls.png}
\caption{Comparison-star control light curves.  Each panel treats one AAVSO
comparison star as the target and divides it by the remaining comparison-star
ensemble.  The shaded band marks the robust scatter for that pseudo-target.
The comparison stars are much more consistent with constant flux than the AM
CVn target light curve, supporting the use of the ensemble for differential
photometry.}
\label{fig:comparison-controls}
\end{figure}

\section{Analysis Methods}

We first fit a model ladder.  The baseline model is constant flux.  The
generic models are a free single sine and a free double-wave model containing a
fundamental plus its first harmonic.  We also ran multi-harmonic
Lomb-Scargle periodograms: \texttt{nterms=2} is the periodogram analogue of the
free double-wave model.  The physically motivated models use literature
periods: the 1028.7322 s orbital period, the 1051.2 s positive superhump, a
1011.4 s negative superhump, and a superhump-plus-harmonic template.  Finally,
we fit a beat-constrained precessing-disc proxy with
\(f_{\rm sh}=f_{\rm orb}-f_{\rm prec}\).  The 13.38 h disc precession period
was not directly fit from the light curve because our run covers only 0.14
cycle of it.

For each model we report chi-square, BIC, RMS residual scatter, and
time-blocked cross-validation error.  For free periods, we compute profile
\(\chi^2\) intervals and fixed-time Monte Carlo perturbation uncertainties.
Formal photon-noise errors underestimate the residual scatter, so the final
Monte Carlo period uncertainties perturb the light curve with the best
double-wave model's \(\sqrt{\chi^2_\nu}=1.59\) empirical error scale.  We
also test stability by fitting the first and second halves separately, clipping
outliers, dropping comparison stars one at a time, and checking residual
periodograms.  As final stress tests, we fit extra nuisance terms for time,
tracking shifts, exposure state, and ensemble common mode; repeat the
double-wave fit with robust loss functions; and rebuild the differential light
curve from comparison-star subsets.

For readers unfamiliar with the model names, all of the plotted curves are
weighted Fourier fits.  For a chosen period \(P_j\) and harmonic number \(h_j\),
we fit
\[
m(t_i)=c+\sum_{j=1}^{M}\left[
a_j\sin\left(2\pi h_jt_i/P_j\right)+
b_j\cos\left(2\pi h_jt_i/P_j\right)\right],
\]
where \(t_i\) is time from the first exposure and \(y_i\) is normalized
differential flux.  The best coefficients minimize
\[
\chi^2=\sum_{i=1}^{N}\left[\frac{y_i-m(t_i)}{\sigma_i}\right]^2 .
\]
A single sine has one sinusoidal term.  The fixed harmonic superhump template
uses the literature superhump period \(P_{\rm sh}=1051.2\) s from Skillman et
al. \cite{skillman1999}, then adds a second sinusoid at twice that frequency.
A free double-wave model has the same mathematical shape---a sinusoid plus
another sinusoid at double its frequency---but it does not fix \(P\) to the
literature value.  The beat-constrained disk model is different in where the
subtraction happens: it does not subtract two brightness sine waves.  Instead,
it fixes the spectroscopic orbital period and scans a precession period through
\[
f_{\rm sh}=f_{\rm orb}-f_{\rm prec}.
\]
This frequency subtraction chooses the superhump period used in the same
double-wave brightness model.  The fitted \(P_{\rm prec}\) is therefore a
conditional beat-relation parameter, not a direct detection of disk precession.

We compare models with an effective BIC,
\[
{\rm BIC}_{\rm eff}=\chi^2+k_{\rm eff}\ln N,\qquad
k_{\rm eff}=k_{\rm lin}+k_{\rm nonlin}.
\]
Here \(k_{\rm lin}=1+2M\) counts the fitted constant, sine, and cosine
coefficients, while \(k_{\rm nonlin}\) counts free scanned quantities such as a
period or precession period.  Lower BIC is better; a model can reduce
\(\chi^2\) and still lose if it spends too many parameters.

\section{Results}

The single-sine search recovers 525.06 s, with fixed-time Monte Carlo
perturbations giving a median of \(524.38^{+5.75}_{-3.84}\) s after empirical
error scaling.  Ordinary Lomb-Scargle detects the same harmonic at 525.08 s.
When Lomb-Scargle is generalized to a two-term Fourier model, it recovers a
fundamental near 1053.27 s.  Thus Lomb-Scargle is useful as a detection tool,
but it must be made multi-harmonic or followed by physical model fitting for AM
CVn.

The best independent model is the free double-wave fit:
\[
P_{\rm free}=1055.21^{+13.10}_{-9.63}~{\rm s}.
\]
This quotes the best-fit period from the original light curve with the
asymmetric fixed-time Monte Carlo uncertainty.  It is consistent with the known
1051.2 s superhump family, but it is not precise enough to improve the
published photometric or spectroscopic periods.

\begin{table}[H]
\centering
\caption{AM CVn model ladder using effective BIC.  \(\Delta{\rm BIC}_{\rm eff}\)
is measured relative to the best free double-wave model; positive values are
worse.}
\label{tab:model-ladder}
\resizebox{\textwidth}{!}{%
\begin{tabular}{llrrrrr}
\toprule
Model & Physical interpretation & Period (s) & \(k_{\rm eff}\) & RMS (\%) & \({\rm BIC}_{\rm eff}\) & \(\Delta{\rm BIC}_{\rm eff}\) \\
\midrule
constant flux & no periodic variability & -- & 1 & 1.66 & 462.78 & 96.75 \\
free single sine & harmonic-only detection & 525.06 & 4 & 1.52 & 403.47 & 37.44 \\
1051.2 s + harmonic & positive-superhump template & 1051.2 & 5 & 1.43 & 361.65 & -4.37 \\
beat-constrained disk & \(f_{\rm sh}=f_{\rm orb}-f_{\rm prec}\) & 1055.22 & 6 & 1.43 & 366.03 & 0.00 \\
free double wave & fundamental + harmonic & 1055.21 & 6 & 1.43 & 366.03 & 0.00 \\
\bottomrule
\end{tabular}
}
\end{table}

The most important result in Table~\ref{tab:model-ladder} is that the
fixed 1051.2 s superhump-plus-harmonic template has the lowest effective BIC.
It earns this advantage because its period is supplied by the literature, not
searched from our data.  The strongest independent period measurement remains
the free double-wave model.  The beat-constrained precessing-disc proxy and the
free double-wave model are effectively tied because the scanned precession
period maps the known orbital period onto nearly the same photometric period;
the derived value is \(P_{\rm prec}\simeq 11.39\) h.  This captures the expected
superhump physics better than a generic Fourier label, but it is not an
independent precession detection because the baseline covers far less than one
precession cycle.

The reduced chi-squared values are not formally good even for the preferred
models.  The constant model has \(\chi^2_\nu=3.34\), the free single sine has
\(\chi^2_\nu=2.84\), and both the free double wave and fixed 1051.2 s harmonic
template have \(\chi^2_\nu\simeq2.53\).  We therefore interpret the periodic
fits as detections of coherent AM CVn-like variability, not as complete
Gaussian-noise descriptions of the light curve.  The residual excess scatter is
why the period uncertainties are empirically scaled and why the final claim is
conservative.

\begin{figure}[H]
\centering
\includegraphics[width=\textwidth]{analysis/am_cvn/plots/AM_CVn_stacked_model_fits.png}
\caption{Light curve with the main model families overplotted.  The constant
and single-sine models leave coherent structure.  The fixed 1051.2 s harmonic
template, beat-constrained disk model, and free double-wave model look very
similar because all describe the same superhump-plus-harmonic waveform over a
111 min baseline.}
\label{fig:model-suite}
\end{figure}

The validation tests are mixed.  Dropping one comparison star at a time changes
the free double-wave period by only 5.2 s, and clipping \(|z|>4\) outliers
changes it by only 0.02 s.  The residual periodogram has false-alarm
probability 0.99, so no strong unmodeled short-period signal remains.  However,
the first-half and second-half fitted periods differ by 22.8 s.  That split is
not surprising with only about six superhump cycles, but it blocks any strong
claim that the period is tightly constrained by our data alone.

We also tried three ways to squeeze more information out of the same light
curve.  First, adding nuisance regressors for time, image shifts, exposure
time, number of comparison stars, and common-mode behavior can reduce the raw
RMS from 1.43\% to 1.41\%, but BIC and blocked cross-validation get worse
(\(\Delta{\rm BIC}=22.66\), \(\Delta{\rm CV}=0.97\%\) for the most flexible
case).  This indicates overfitting rather than a better physical constraint.
Second, robust likelihoods do not sharpen the answer: soft-\(L_1\), Huber, and
Cauchy losses shift the best period to 1057--1059 s and increase the ordinary
RMS.  Third, selected comparison-star subsets can shave the target RMS as low
as 1.40\%, but the fitted period moves by a few seconds and the control-star
scatter does not justify replacing the full ensemble.  These tests are useful
because they show what does not work: the limiting uncertainty is the short
baseline and cycle-to-cycle structure, not just Gaussian photometric noise.

The broadband photometry cannot test AM CVn's S-wave emission directly.  The
S-wave is a spectroscopic velocity feature from the bright spot or stream-disc
impact region.  Testing it requires time-resolved spectroscopy or Doppler
tomography, not V-band imaging.  For the same reason, our photometry can support
the known white-dwarf-binary interpretation but cannot independently prove or
disprove the object class.

\begin{figure}[H]
\centering
\includegraphics[width=0.78\textwidth]{analysis/am_cvn/plots/AM_CVn_physical_literature_constraints.png}
\caption{Comparison between our period constraints and literature periods.  Our
independent Monte Carlo uncertainty is about 10--13 s, far larger than the
published precision of the orbital and superhump periods.  The result is
therefore consistent with the literature, not more constraining than it.}
\label{fig:literature-constraints}
\end{figure}

\section{Conclusions}

Our data clearly detect short-timescale AM CVn variability.  A single-sine
model finds a 525 s signal, while the physically preferred interpretation is a
double-wave signal in the known 1051 s superhump family.  The best independent
period is \(1055.21^{+13.10}_{-9.63}\) s, and the literature-fixed 1051.2 s
superhump-plus-harmonic template is the best conditional model by effective BIC.

The result does not tighten the literature constraints.  The physically
motivated fixed-period model is a better conditional description, but the
first/second-half stability test fails and our independent uncertainty is
hundreds of times larger than the best published period precision.  Detrending,
robust likelihoods, and comparison-star subset searches improve neither
predictive performance nor period stability.  The strongest defensible claim is
therefore: our photometry independently recovers variability consistent with AM
CVn's known superhump/harmonic structure, but a longer multi-night baseline
would be needed to improve published period constraints.

\section*{Acknowledgments}

We thank the PHYSICS 100 teaching staff for telescope access and reduction
guidance.  This work used AAVSO comparison-star resources and NASA ADS/arXiv
literature records.

\begin{thebibliography}{9}

\bibitem{nelemans2001}
Nelemans, G., Steeghs, D., \& Groot, P. J. 2001.
\textit{Spectroscopic evidence for the binary nature of AM CVn}.
MNRAS, 326, 621.
\url{https://academic.oup.com/mnras/article/326/2/621/1006901}.

\bibitem{patterson1979}
Patterson, J., Nather, R. E., Robinson, E. L., \& Handler, F. 1979.
\textit{Increase in the orbital period of AM Canum Venaticorum}.
ApJ, 232, 819.
\url{https://articles.adsabs.harvard.edu/pdf/1979ApJ...232..819P}.

\bibitem{roelofs2006}
Roelofs, G. H. A., Groot, P. J., Nelemans, G., Marsh, T. R., \& Steeghs, D. 2006.
\textit{Kinematics of the ultracompact helium accretor AM Canum Venaticorum}.
MNRAS, 371, 1231.
\url{https://academic.oup.com/mnras/article/371/3/1231/1007944}.

\bibitem{skillman1999}
Skillman, D. R., Patterson, J., Kemp, J., Harvey, D. A., Fried, R. E.,
Retter, A., Lipkin, Y., \& Vanmunster, T. 1999.
\textit{Superhumps in Cataclysmic Binaries. XVII. AM Canum Venaticorum}.
PASP, 111, 1281.

\bibitem{smak2023}
Smak, J. 2023.
\textit{AM CVn -- System Parameters and Gravitational Waves}.
Acta Astronomica, 73, 1.
\url{https://arxiv.org/pdf/2312.16506}.

\bibitem{solheim2010}
Solheim, J.-E. 2010.
\textit{AM CVn Stars: Status and Challenges}.
PASP, 122, 1133.

\end{thebibliography}

\appendix
\section{Individual Contributions}

TODO: Fill in team-member names and individual contributions before
submission, as required by the project instructions.

\section{Reproducibility}

The model suite is generated with:
\begin{verbatim}
PYTHONPATH=src python -m amcvn.sequence_analysis
PYTHONPATH=src python -m amcvn.physical_model_suite --bootstrap 250 --monte-carlo 250
PYTHONPATH=src python -m amcvn.stacked_model_fits
\end{verbatim}
The generated CSV files use the \texttt{physical\_*} prefix in the AM CVn
periodogram folder, including the primary
\texttt{physical\_monte\_carlo\_period\_summary.csv}.  The generated figures
use the \texttt{AM\_CVn\_*} prefix in the plot folder.

\end{document}
